% Proposed after checkpoint10; authored in work_CE.
\statement{A.70}{Proved}{Finite actual central data do not suffice for unrestricted equivariant peripheral lifting}
Let
\[
 \mathcal E=F\rtimes\langle c\rangle
   \cong C_2\amalg C_2\amalg C_2,
 \qquad N=\ker\bigl(F\longrightarrow(\mathbb Z/12)^2\bigr).
\]
There is an involutory automorphism of the actual source-central
quotient \(\mathcal E/[N,\mathcal E]\) which fixes the external
reflection, preserves all three named reflection classes, and
preserves the three named pair-product cusp classes with one common
unit multiplier. It induces an automorphism of \(\mathcal E/N\),
but has no full reflection-class-preserving peripheral lift to
\(\mathcal E\), even if the full lift is allowed infinite order.
The multiplier is mixed: its residue is \(7\pmod {24}\).
This does not refute ordinary R2, or a multiplier-one equivariant
rule for one fixed action. It makes the equivariance failure of the
mixed-sign construction in A.18 explicit on a finite actual central
refinement, using the fixed-group argument of A.6 and the S4 test
of A.39.

\begin{proof}
Put \(w=[x,y]=xyx^{-1}y^{-1}\). Since \(F'\subseteq N\),
killing \([N,F]\) makes the orientation subgroup class two.
It also imposes \(w^{12}=1\). Conversely, in the class-two group
with this relation, \(x^{12},y^{12},w\) are central and generate
the image of N, so these are exactly the required relations.
The action of c inverts \(x^{12},y^{12}\) and fixes w. Hence
killing \([N,c]\) adds \(x^{24}=y^{24}=1\) and no other
relation. The exact quotient has normal forms
\[
 x^a y^b w^d c^e,
 \qquad a,b\in\mathbb Z/24,\quad d\in\mathbb Z/12,\quad
 e\in\mathbb Z/2,
 \tag{A.70.1}
\]
with multiplication
\[
 (a,b,d,e)(u,v,k,f)=
 (a+(-1)^e u,\ b+(-1)^e v,\ d+k-b(-1)^e u,\ e+f).
 \tag{A.70.2}
\]
This explicit group realizes the converse quotient, so there are
no unproved additional relations. In particular
\[
 |\mathcal E/[N,\mathcal E]|=2\cdot24^2\cdot12=13824,
 \qquad N/[N,\mathcal E]\cong C_2^2\times C_{12}.
 \tag{A.70.3}
\]
The central kernel is generated by \(x^{12},y^{12},w\).

Define
\[
 \theta(x)=x^7,\qquad\theta(y)=y^7,\qquad\theta(c)=c.
 \tag{A.70.4}
\]
In (A.70.1) this is \((a,b,d,e)\mapsto(7a,7b,d,e)\),
because \(49\equiv1\pmod {12}\). It is an automorphism of
order two, since \(49\equiv1\pmod {24}\), and fixes the
entire central kernel in (A.70.3). Thus it really lies over an
automorphism of the original finite quotient
\((\mathbb Z/12)^2\rtimes\langle c\rangle\).
The three named reflection classes may be represented by
\(c,xc,y^{-1}c\); their images under \(\theta\) are their conjugates
by \(1,x^3,y^{-3}\), respectively. If \(z=(xy)^{-1}\),
the class-two formula gives
\[
 (xy)^7=x^7y^7w^{-21},\qquad
 \theta(xy)=x^3(xy)^7x^{-3},\qquad
 \theta(z)=x^3z^7x^{-3}.
 \tag{A.70.5}
\]
The residue 7 is the image of the genuine full unit
\(\lambda\in\widehat{\mathbb Z}^{\times}\) with
\(\lambda_2=-1\) and \(\lambda_p=1\) for every odd p.
Consequently these are common-unit peripheral data, not a
substitution by a nonunit integer.

Suppose a full lift existed. Preserving the named reflection
classes preserves the orientation subgroup F, because its parity
character is one on each of those classes. The image of c is
conjugate to c, and its conjugator can be chosen in F: multiply
an odd conjugator by c if necessary. Conjugating the full lift
by this inner element makes it fix c, without changing its
peripheral multiplier \(\Lambda\). The orientation
abelianization of (A.70.1) forces
\(\Lambda\equiv7\pmod {24}\); this remains true after the
inner normalization.

Write its x-image as \(a x^\Lambda a^{-1}\), with \(a\in F\).
Commuting with c implies
\(a^{-1}c(a)\in C_F(x)=\overline{\langle x\rangle}\),
say \(a^{-1}c(a)=x^t\). If \(a^{\rm ab}=(r,s)\),
abelianization yields
\[
 (t,0)=(-2r,-2s),\qquad s=0,\quad t=-2r.
\]
Thus \(a x^{-r}\) is literally c-fixed. The established
full identity \(C_F(c)=1\) gives \(a=x^r\). The x-image
is therefore exactly \(x^\Lambda\); the same argument
makes the y-image exactly \(y^\Lambda\). This uses the full
fixed-group theorem and does not assert trivial fixed groups
in finite quotients.

Evaluate at \(X=(2\,3\,4)\), \(Y=(1\,2\,3\,4)\) in
S4, with rightmost permutations acting first. Since
\(\Lambda\equiv7\pmod {12}\), the element
\(X^\Lambda Y^\Lambda\) has order two, whereas
\((XY)^\Lambda\) has order four. They cannot be conjugate,
contradicting the third-cusp condition. No order assumption on
the proposed full lift was used.

The self-contained certificate
\path{work_CE/research_equivariant_R2/check_finite_central_equivariant_counterexample.py}
checks all 13824 normal forms and generator edges, the 48-element
central kernel, the reflection and cusp conjugators, and the S4
obstruction. Its stdout and JSON are stored alongside it. The
presentation argument above identifies the actual source-central
quotient; the finite count alone is not used to infer that identity.

The transported-head formulation is essential. There is an
\emph{original} epimorphism from the actual three-reflection
source to \(\mathcal E/N\), and \(\theta\) transports its named
quotient tuple. The desired map on the entire actual central
refinement already exists. The theorem disproves a rule that
would infer a full equivariant peripheral lift from these data.
It leaves the same-action multiplier-one restriction undecided.
\end{proof}

\paragraph{A multiplier-one variant between different full actions.}
This scope distinction can also be made on actual full involutions.
Choose an arithmetic \(\widehat{GT}\) element
\(\psi=(\lambda,f)\), \(f\in F'\), with the mixed unit
\(\lambda\) used above; the usual arithmetic embedding and cyclotomic
surjectivity supply such an element. In every class-two quotient
\(f\) is central. Therefore \(\psi\) acts on the \emph{entire} orientation
subgroup in (A.70.1) exactly as \(\theta\), not merely on its abelian
quotient or cusp classes. Put \(\alpha=\psi c\psi^{-1}\).
This is a genuine full negative involution. Power substitution
commutes with c in a class-two group, so \(\alpha\) and \(c\) have identical
actions on \(F/\gamma_3(F)\). Since
\(\gamma_3(F)\subseteq[N,F]\), they have identical actions
on \(N/[N,F]\), and hence
\[
 [N,F\rtimes\langle\alpha\rangle]
       =[N,F\rtimes\langle c\rangle].
 \tag{A.70.6}
\]
Their actual source-central quotients are literally identified.
The identity on this quotient is an equivariant multiplier-one
map: it fixes the orientation cusp elements and external
reflection, and the other named reflection classes match by
(A.70.5) and the displayed reflection conjugators.
But a full multiplier-one conjugator g with
\(gcg^{-1}=\alpha\) would make \(\psi^{-1}g\) centralize c
in the full outer peripheral group. By A.6 this centralizer is
\(\{1,c\}\), whose multipliers are \(1,-1\), whereas
\(\lambda^{-1}\notin\{1,-1\}\). Such g cannot exist.
This concerns \emph{two different source actions}; it is not a
counterexample to a multiplier-one lifting rule for the single
fixed action c, and \(\alpha\) is already conjugate to \(c\) by \(\psi\).

\statement{A.71}{Proved}{Three-trace rigidity fixes every residually irreducible rank-two matrix quotient}
Let A be a finite commutative ring with identity and let
\(\rho:F\to\operatorname{SL}_2(A)\) be a continuous
representation. Assume its reduction at every maximal ideal of
A is absolutely irreducible. For every full peripheral
automorphism \(\Phi\) with common multiplier \(+1\) or \(-1\),
there is \(G\in\operatorname{GL}_2(A)\) such that
\[
                   \rho\Phi=\operatorname{Ad}(G)\rho.
 \tag{A.71.1}
\]
In particular \(\ker\rho\) is stable as a conclusion, without
any prior stability or semilinear-action hypothesis. Every
epimorphism \(F\to\operatorname{PSL}_2(k)\), for a finite
field \(k\) of order at least four, has a stable kernel under
such \(\Phi\), and its induced action is projective linear.
The same conjugacy conclusion holds over a profinite commutative
ring when the stated residual hypothesis holds at every finite
quotient. This is an application of classical three-trace
rigidity, with the local-ring and inverse-limit details supplied
below; no novelty claim is made for the trace theorem.
\cite{CEMcCulloughWanderley2013}

\begin{proof}
Set
\[
 X=\rho(x),\quad Y=\rho(y),\quad
 X'=\rho\Phi(x),\quad Y'=\rho\Phi(y).
\]
Conjugation preserves trace, and inversion in \(\operatorname{SL}_2\)
also preserves trace over every commutative ring. The three
peripheral conditions therefore give
\[
 \operatorname{tr}X'=\operatorname{tr}X=a,\qquad
 \operatorname{tr}Y'=\operatorname{tr}Y=b,\qquad
 \operatorname{tr}(X'Y')=\operatorname{tr}(XY)=d.
 \tag{A.71.2}
\]
Because \(\Phi\) is onto, the primed pair generates the same image as
the unprimed pair, and has the same residual irreducibility.
Cayley--Hamilton gives the identities
\[
 X^2=aX-I,\quad Y^2=bY-I,\quad
 YX=-XY+bX+aY+(d-ab)I.
 \tag{A.71.3}
\]
These hold in characteristic two as well. They imply that every
word in X and Y lies in the A-span of \(I,X,Y,XY\), and that
all multiplication in this spanning set depends only on
\(a,b,d\). Modulo any maximal ideal, absolute irreducibility
makes the generated algebra the full matrix algebra. Here is an
elementary two-dimensional verification after extension to an
algebraic closure. If a nonscalar generator has distinct
eigenvalues, its two spectral idempotents belong to the algebra;
irreducibility forces both off-diagonal entries of the other
generator to be nonzero, giving all four matrix units. If instead
it has a repeated eigenvalue, subtracting that eigenvalue and
changing basis gives \(E_{12}\). The lower-left entry of the
other generator is nonzero by irreducibility, and multiplying
by \(E_{12}\) then gives \(E_{11}\), hence the remaining
matrix units. Both generators cannot be scalar. Dimension
descends to the residue field, so the four matrices above are
a basis there.
Their coordinate determinant is consequently outside every
maximal ideal of A, hence is a unit. It follows that
\(I,X,Y,XY\) is an A-basis of \(M_2(A)\). The same holds
for the primed matrices. Equations (A.71.2)--(A.71.3) now define
an A-algebra automorphism
\[
 M_2(A)\longrightarrow M_2(A),\qquad X\mapsto X',\quad Y\mapsto Y'.
\]
This passage is a comparison of complete multiplication tables,
not a claim that traces distinguish arbitrary reducible pairs.

For completeness this algebra automorphism is inner over A.
A finite commutative ring is a product of finite local rings,
and an A-algebra map fixes the central idempotents, so work in
one local factor. The images \(P_1,P_2\) of the standard
matrix units \(E_{11},E_{22}\) are complementary idempotents.
Their image modules in \(A^2\) are rank-one projective
modules, hence free over the local ring. Choose a generator
\(v_1\) of the first and put \(v_2=\varphi(E_{21})v_1\).
The matrix-unit identities show that \(v_2\) generates the
second and \(\varphi(E_{12})v_2=v_1\). The matrix with columns
\(v_1,v_2\) is invertible and implements phi on every matrix
unit. Combining these matrices over the local factors gives
G in (A.71.1). Equality of the representations follows from
equality on the free generators.

For a projectively onto representation over k, choose arbitrary
determinant-one lifts of the two free generator images. Their
projective image is \(\operatorname{PSL}_2(k)\), so their
matrix action has no common invariant line over an algebraic
closure. The finite-field result applies to this lifted
representation, and projectivizing proves both stability of the
original kernel and projective linearity. It is unnecessary to
assert that these particular lifts generate all of
\(\operatorname{SL}_2(k)\).

Finally let A be an inverse limit of its finite quotient rings
\(A_i\), and apply the finite result to \(\rho_i\). Let
\[
 \mathcal C_i=\{G_i\in\operatorname{GL}_2(A_i):
        \rho_i\Phi=\operatorname{Ad}(G_i)\rho_i\}.
\]
Each is a nonempty finite set. Reduction maps carry
\(\mathcal C_j\) to \(\mathcal C_i\); surjectivity of
these maps is not assumed. Any finite list of compatibility
conditions is satisfied by choosing a conjugator at one common
finite refinement and reducing it. Compactness of the product
therefore gives a compatible family. Its inverse family of
matrices is compatible as well, so the limit lies in
\(\operatorname{GL}_2(A)\) and gives the required conjugacy.
This argument supplies a single conjugator in the inverse limit,
not merely unrelated conjugators at each finite level.

Over finite fields the classical nonsingular trace-triple
conjugacy result is recorded as Macbeath's theorem in
McCullough--Wanderley, Theorem 8.2, pp.~490--491.
\cite{CEMcCulloughWanderley2013}
The elementary algebra proof above gives the precise
\(\operatorname{GL}_2\) conclusion and its finite-ring
extension; it does not incorrectly assert a unique
\(\operatorname{SL}_2(k)\)-conjugacy class in odd characteristic.
Nothing here applies to the canonical unmarked substitution
before its third-cusp condition is proved, or to arbitrary
higher-dimensional finite-group quotients.
\end{proof}
